\section{从辽东到燕云：战争怎样把道路变成边界}
\label{sec:frontier-war}

辽水以东的城防、河流和疫病，比诏书上的疆界更早决定了隋军能够走多远。598年的水陆出兵已因风浪、疾病和补给撤回；612年大举东征，辽东城久攻不下，渡水部队覆败，613年黎阳的政变又迫使朝廷把目光从前线转回腹地。高句丽的城防与纵深并非被动背景，隋廷也不是面对一个可以任选的战场：它若继续追求速胜，便要把更多兵粮压到漫长道路上；若撤回，则承认边境竞争没有被统一帝国的名义消除。%
\footnote{隋与高句丽的三次关键行动见 ST-598-GOGURYEO-FIRST-CAMPAIGN、ST-612-GOGURYEO-INVASION、ST-614-GOGURYEO-THIRD，依据《隋书》本纪及〈东夷传·高丽〉、《资治通鉴》。高句丽方面连续叙事缺环，兵数、船数与伤亡不得按隋方战史实数化。}

唐初的边疆收获也须放在道路上读。630年破东突厥、640年置西州、657年击西突厥，使唐廷得以在草原与天山南北安排不同深度的安置、都护和州县；每一种安排都要支付粮草、任官和交涉的成本。668年唐与新罗攻破平壤后设置安东都护府，恰说明战场胜利不是接收的终点：唐罗之间立即围绕旧地的控制发生新的摩擦。唐方本纪和列传擅长记录册封与战功，新罗材料、吐鲁番文书、碑铭与城址则提醒人们，都护府名称不能替代地方社会的服从。%
\footnote{东突厥、高昌、西突厥和高句丽覆亡见 ST-630-EASTERN-TURKS、ST-640-GAOCHANG、ST-657-WESTERN-TURKS、ST-668-GOGURYEO-CONQUEST。可对读两《唐书》相关本纪、外族传与《三国史记》；安置人口、有效辖区和盟军贡献存在材料差异。}

751年的怛罗斯之败与云南方向的挫折，表明唐的军令已经越过了本国能够稳定供给的尺度；但它们并不自动预告755年安禄山的起兵。安史危机来自另一种更近的边疆化：范阳、平卢、河东三镇的兵员、仓廪、将校网络集中在同一指挥者周围，朝廷在东北防务、宫廷人事和关中供给之间失去足够的回旋空间。潼关失守后，灵武朝廷必须借朔方军、地方守备与回纥援军再造作战线；763年吐蕃进入长安，又表明收复两京并未把西北道路重新锁回京师。%
\footnote{怛罗斯、安禄山起兵、潼关、灵武和吐蕃入长安见 ST-751-TALAS、ST-755-AN-LUSHAN-REBELLION、LT-756-001、LT-756-002、LT-763-014。核心为两《唐书》本纪、安禄山与哥舒翰诸传、《资治通鉴》，并应与突厥、回纥、吐蕃碑铭和中亚材料相互校正。}

战后唐廷并未简单地从边地退出。它以节度使名义同河北旧将议价，又在河西等待新的地方中介。848年张议潮逐出沙州吐蕃守军，851年才获授归义军节度使；这段距离和时间说明册命能重新连结道路，却不足以恢复常规州县直辖。五代时，燕云十六州进入石敬瑭换取契丹援军的交易，问题又从河西走廊转到开封北面的城、牧地、互市与军费。把这类选择只写作屈辱或背叛，会遗漏当事人首先面对的兵力和供给约束；把它们写作单纯的现实主义，也会遮住边地居民承担的驻防、迁徙和战争代价。%
\footnote{归义军与燕云问题见 LT-848-051、LT-851-052、LT-936-099。《旧唐书·张议潮传》《唐会要》卷78〈节度使〉、《旧五代史·晋高祖纪一》《辽史·太宗纪上》与《资治通鉴》分别保存这些行动；“十一州”、燕云范围和双方礼仪措辞均不能脱离其政治语境。}

\paragraph{边疆不是帝国边缘的空白。}
战争带来的制度收益，是迫使国家学会利用道路、仓储、盟军与地方中介；唐的都护、节度与册命因此能在远处维持联系。它的损失同样制度化：远征会抽走腹地兵粮，边军会积累独立的供给和任用能力，暂时的援军或割地协议会改写下一轮战争的起点。辽东、河西、范阳和燕云并非同一种“边患”，却都说明版图上的颜色不等于能够持续支付的控制。
